\documentclass[11pt]{article}

\usepackage[utf8]{inputenc}
\usepackage[T1]{fontenc}
\usepackage{lmodern}
\usepackage{geometry}
\usepackage{amsmath,amssymb,amsfonts,amsthm,mathtools}
\usepackage{bm}
\usepackage{enumitem}
\usepackage{microtype}
\usepackage{hyperref}
\usepackage{titlesec}
\usepackage{booktabs}
\usepackage{array}
\usepackage{setspace}
\usepackage{mathrsfs}
\usepackage{physics}

\geometry{margin=1in}
\hypersetup{colorlinks=true, linkcolor=black, urlcolor=blue, citecolor=black}
\setlist[enumerate]{topsep=0.4em,itemsep=0.35em}
\setlist[itemize]{topsep=0.3em,itemsep=0.25em}
\titleformat{\section}{\large\bfseries}{\thesection.}{0.5em}{}
\titleformat{\subsection}{\normalsize\bfseries}{\thesubsection.}{0.5em}{}
\setstretch{1.06}

\newcommand{\Aether}{\AE{}ther}
\newcommand{\Aflow}{\AE{}ther-flow}
\newcommand{\TheoryName}{\Aflow{} Interpretation of Relativity}
\newcommand{\TheoryProgram}{\Aether{} / \Aflow{} framework}

\newcommand{\CompletedMetric}{g^{\sharp}}
\newcommand{\MatterSpace}{\Psi}

\newcommand{\Car}{\mathrm{car}}
\newcommand{\Cur}{\mathrm{cur}}
\newcommand{\Hom}{\mathrm{hom}}

\newcommand{\ForSpace}[1]{\mathcal I_{#1}^{\mathrm{for}}}
\newcommand{\PrimShell}{\mathcal S_{\mathrm{prim}}}
\newcommand{\Reduction}[1]{\mathcal R_{#1}}
\newcommand{\CurrentCarrier}[1]{\mathfrak C_{#1}^{\mathrm{cur}}}
\newcommand{\EqCarrier}[1]{\mathfrak C_{#1}^{\mathrm{eq}}}
\newcommand{\MetricOut}{\mathscr G_{\mathrm{out}}}
\newcommand{\MatterOut}{\mathscr T_{\mathrm{out}}}

\newcommand{\SubTarget}[1]{E_{#1}^{\mathrm{sub}}}
\newcommand{\SubPair}[1]{\mathfrak s_{#1}^{\mathrm{sub}}}
\newcommand{\EqnTarget}[1]{Q_{#1}^{\mathrm{eqn}}}
\newcommand{\HiddenSpace}[1]{\mathcal H_{#1}^{\mathrm{hid}}}
\newcommand{\HiddenBody}[1]{D_{#1}^{\mathrm{hid}}}

\newcommand{\ProtoSpace}[1]{\mathcal U_{#1}^{\mathrm{proto}}}
\newcommand{\ProtoBody}[1]{\mathcal B_{#1}^{\mathrm{proto}}}
\newcommand{\ProtoSection}[1]{\Gamma_{#1}^{\mathrm{proto}}}
\newcommand{\ProtoShellSet}[1]{\mathcal Z_{#1}^{\mathrm{proto}}}
\newcommand{\ProtoZeroCore}[1]{\mathcal Z_{#1}^{0,\mathrm{proto}}}
\newcommand{\ProtoSplit}[1]{\Phi_{#1}^{\mathrm{proto}}}
\newcommand{\ProtoCharge}[1]{\mathcal Q_{#1}^{\mathrm{proto}}}
\newcommand{\ProtoEmbed}[1]{\zeta_{#1}^{\mathrm{proto}}}
\newcommand{\ProtoKernelCh}[1]{\widetilde K_{#1}^{0,\mathrm{proto}}}
\newcommand{\ProtoStateComp}[1]{P_{#1}^{\mathrm{proto,co}}}

\newcommand{\PreProtoNucleus}{\mathfrak Q^{\mathrm{nuc}}}
\newcommand{\PreProtoSpace}[1]{\mathcal V_{#1}^{\mathrm{preproto}}}
\newcommand{\PreProtoBody}[1]{\mathcal B_{#1}^{\mathrm{preproto}}}
\newcommand{\PreProtoProj}[1]{\Pi_{#1}^{\mathrm{preproto}}}
\newcommand{\PreProtoProtoMin}[1]{\mathfrak f_{#1}^{\mathrm{preproto,proto}}}
\newcommand{\PreProtoSection}[1]{\Gamma_{#1}^{\mathrm{preproto}}}
\newcommand{\PreProtoFunctional}[1]{\mathscr W_{#1}^{\mathrm{preproto}}}
\newcommand{\PreProtoShellSet}[1]{\mathcal Z_{#1}^{\mathrm{preproto}}}

\newcommand{\PreNucCore}{\mathfrak R^{\mathrm{core}}}
\newcommand{\PreNucSpace}[1]{\mathcal W_{#1}^{\mathrm{prenuc}}}
\newcommand{\PreNucBody}[1]{\mathcal B_{#1}^{\mathrm{prenuc}}}
\newcommand{\PreNucProj}[1]{\Pi_{#1}^{\mathrm{prenuc}}}
\newcommand{\PreNucNucleusMin}[1]{\mathfrak f_{#1}^{\mathrm{prenuc,nuc}}}
\newcommand{\PreNucShellSet}[1]{\mathcal Z_{#1}^{\mathrm{prenuc}}}

\newcommand{\PreCoreTransportCore}{\mathfrak S^{\mathrm{tr}}}
\newcommand{\PreCoreTrBody}[1]{\mathcal B_{#1}^{\mathrm{precore,tr}}}
\newcommand{\PreCoreTrProj}[1]{\widehat{\Pi}_{#1}^{\mathrm{precore}}}
\newcommand{\PreCoreTrLift}[1]{\mathfrak f_{#1}^{\mathrm{precore,tr}}}
\newcommand{\PreCoreTrShell}[1]{\mathcal Z_{#1}^{\mathrm{precore,tr}}}

\newcommand{\SubPreSectionCore}{\mathfrak U^{\mathrm{sec}}}
\newcommand{\SubPreSection}[1]{\mathfrak s_{#1}^{\mathrm{subpre,sec}}}

\newcommand{\SubSectionPackage}{\mathfrak V^{\mathrm{subsec}}}
\newcommand{\SubSectionSpace}[1]{\mathcal M_{#1}^{\mathrm{subsec}}}
\newcommand{\SubSectionBody}[1]{\mathcal B_{#1}^{\mathrm{subsec}}}
\newcommand{\SubSectionProj}[1]{\Pi_{#1}^{\mathrm{subsec}}}
\newcommand{\SubSectionFunctional}[1]{\mathscr E_{#1}^{\mathrm{subsec}}}
\newcommand{\SubSectionFiber}[2]{\mathcal F_{#1}^{\mathrm{subsec}}(#2)}
\newcommand{\SubSectionShell}[1]{\mathcal Z_{#1}^{\mathrm{subsec}}}

\newtheorem{definition}{Definition}
\newtheorem{proposition}{Proposition}
\newtheorem{remark}{Remark}
\newtheorem{corollary}{Corollary}
\newtheorem{theorem}{Theorem}

\title{\textbf{The \TheoryName{}}\\[0.4em]
\large Primitive-Reservoir Observer-Reduction\\
\large Observer-Relevant Sub-Pre-Core Exact-Shell Transport-Section Core\\
\large from Exact-Shell-Transport-Section-Generating Substrate Sub-Section Dynamics}
\author{}
\date{}

\begin{document}

\maketitle

\begin{abstract}
The current primitive-reservoir observer-reduction frontier already sharpened the live burden inside the exact-shell-transport-core-generating substrate sub-pre-core line: the separate transport-image presentation was shown to collapse to the smaller observer-relevant sub-pre-core exact-shell transport-section core. After that theorem, another same-output deeper-origin relay that merely preserves the same section core no longer counts as the honest next benchmark-facing move. The next admissible task is to derive or justify that section core itself from more primitive substrate structure, or else prove a still smaller collapse strictly inside it.

The present note takes the direct derivation route. For each sector it introduces an exact-shell-transport-section-generating substrate sub-section dynamics package consisting of a compact convex body over the already recorded pre-core exact-shell transport body, an affine projection to that fixed body, a nonnegative smooth fiber-selection functional with unique fiberwise minimizers, and an exact shell locus given precisely by the minimizing transport shell. The current section core is then no longer an independent primitive stopping point. It is the canonical minimizer section selected by that deeper fiberwise substrate law.

The benchmark boundary remains fixed throughout. The one operative metric is still exactly \(\CompletedMetric\), the ordinary matter route is unchanged, there is no second operative metric, the low-energy matter law is not enlarged, and no stronger promotion language is licensed. The positive result is therefore a genuine benchmark-facing gain below the current section-core frontier: the remaining open burden is no longer the observer-relevant sub-pre-core exact-shell transport-section core as such, but the deeper exact-shell-transport-section-generating substrate sub-section dynamics that canonically generate it.
\end{abstract}

\section{Introduction}

The active derivational program of The \TheoryName{} has again reached a sharper burden than another same-output relay. The current active-primary theorem already removed the separate transport-image presentation as the live benchmark-facing datum by showing that, once the recorded pre-core exact-shell transport body and shell are fixed, the only independent observer-relevant datum left at that level is the canonical transport section itself. That theorem matters precisely because it blocks a familiar form of drift: depth alone no longer counts once the live section-core datum is unchanged.

The next honest move is therefore not to relay the old image core through a new ambient package. It is to explain why the current section core itself should hold. The present note makes that move in the most direct form now available without changing the benchmark claim boundary. It introduces a deeper substrate package whose defining datum is not another prelisted section, image body, or restricted projection. Instead, each point of the fixed pre-core transport body carries a compact selector fiber and a nonnegative fiber-selection functional. Strict convexity on each fiber forces a unique minimizer, and that minimizing locus is exactly the current transport section.

\paragraph{Framework and claim boundary.}
\TheoryProgram{} denotes the broader \Aether{} / \Aflow{} framework built on the claims that the \Aether{} is the underlying four-dimensional substrate of reality, the \Aflow{} is its intrinsic ordered motion, observed three-dimensional space is the local experiential slice of that deeper substrate, S-time is the experienced order of change arising from matter, light, and the \Aflow{}, and observed expansion is the three-dimensional appearance of deeper four-dimensional ordered motion. Within that framework, \TheoryName{} denotes the exact-closure theory in which the effective gravitational dynamics are adopted to be exactly Einsteinian with universal matter coupling. In the active sequence, \emph{adoption} means use of that established relativistic dynamics without claiming substrate derivation, while \emph{derivation} is reserved for a first-principles recovery from explicit substrate variables.

The present note stays strictly inside that boundary. It does not introduce a second operative metric, it does not enlarge the low-energy matter law, and it does not reopen the already-screened larger image-core, sub-pre-core, or pre-core packages as if they were again the live burden. Its narrower positive role is to show that the current transport section is itself canonically generated by a deeper fiberwise substrate selection law.

\section{Fixed Inputs}

\begin{definition}[Recorded primitive continuation data]
\label{def:fixed}
Fix the following continuation data from the active primitive-reservoir observer-reduction line.
\begin{enumerate}
    \item For each sector label \(\sigma\in\{\Car,\Cur,\Hom\}\), a primitive forcing shell
    \[
    \ForSpace{\sigma},
    \]
    a visible reduction map
    \[
    \Reduction{\sigma}:\ForSpace{\sigma}\to X_{\sigma},
    \]
    and shell-level current and equilibrium carriers
    \[
    \CurrentCarrier{\sigma}:\ForSpace{\sigma}\to Z_{\sigma}^{\mathrm{cur}},
    \qquad
    \EqCarrier{\sigma}:\ForSpace{\sigma}\to Z_{\sigma}^{\mathrm{eq}}.
    \]
    \item The product shell
    \[
    \PrimShell
    \equiv
    \ForSpace{\Car}\times \ForSpace{\Cur}\times \ForSpace{\Hom}.
    \]
    \item The recorded metric and ordinary-matter outputs
    \[
    \MetricOut:\PrimShell\to\{\text{observer metrics}\},
    \qquad
    \MatterOut:\PrimShell\times\MatterSpace\to\{\text{observer stress tensors}\},
    \]
    satisfying
    \begin{equation}
    \MetricOut(\mathbf w)=\CompletedMetric,
    \qquad
    \MatterOut(\mathbf w,\psi)
    =
    T_{\mu\nu}^{\mathrm{matter}}[\CompletedMetric,\psi]
    \label{eq:fixedoutputs}
    \end{equation}
    for every \(\mathbf w\in\PrimShell\) and every matter configuration \(\psi\in\MatterSpace\).
\end{enumerate}
\end{definition}

\begin{definition}[Recorded proto-germ exact-shell target]
\label{def:proto}
Fix \(\sigma\in\{\Car,\Cur,\Hom\}\). The active line already fixes:
\begin{enumerate}
    \item the same substrate target and pairing
    \[
    \SubTarget{\sigma},
    \qquad
    \SubPair{\sigma}:
    \SubTarget{\sigma}\times\SubTarget{\sigma}\to\mathbb R;
    \]
    \item a hidden fiber space \(\HiddenSpace{\sigma}\) with compact hidden body \(\HiddenBody{\sigma}\subset\HiddenSpace{\sigma}\);
    \item a finite-dimensional proto-germ state space \(\ProtoSpace{\sigma}\), a compact convex proto-germ body \(\ProtoBody{\sigma}\subset\ProtoSpace{\sigma}\), a smooth proto-germ restriction section
    \[
    \ProtoSection{\sigma}:\ProtoSpace{\sigma}\to\SubTarget{\sigma},
    \]
    and the exact proto-germ shell
    \[
    \ProtoShellSet{\sigma}
    \equiv
    \bigl(\ProtoSection{\sigma}\bigr)^{-1}(0)\cap\ProtoBody{\sigma};
    \]
    \item a closed embedded proto zero core \(\ProtoZeroCore{\sigma}\subset\ProtoShellSet{\sigma}\) and a split diffeomorphism
    \[
    \ProtoSplit{\sigma}:
    \ProtoZeroCore{\sigma}\times\HiddenBody{\sigma}
    \xrightarrow{\cong}
    \ProtoShellSet{\sigma};
    \]
    \item a hidden charge
    \[
    \ProtoCharge{\sigma}:\HiddenSpace{\sigma}\to\EqnTarget{\sigma},
    \]
    a smooth observer-compatible exact proto-germ zero-response realization
    \[
    \ProtoEmbed{\sigma}:\ForSpace{\sigma}\hookrightarrow\ProtoZeroCore{\sigma},
    \]
    and a fixed-data solvability / coercive package on \(\ProtoZeroCore{\sigma}\), recorded through \(\ProtoKernelCh{\sigma}\) and \(\ProtoStateComp{\sigma}\).
\end{enumerate}
\end{definition}

\begin{definition}[Recorded current transport nucleus]
\label{def:nucleus}
Fix \(\sigma\in\{\Car,\Cur,\Hom\}\). The recorded current transport nucleus \(\PreProtoNucleus\) for sector \(\sigma\) consists of:
\begin{enumerate}
    \item the pre-proto-germ state space \(\PreProtoSpace{\sigma}\), compact convex body \(\PreProtoBody{\sigma}\subset\PreProtoSpace{\sigma}\), projection
    \[
    \PreProtoProj{\sigma}:\PreProtoSpace{\sigma}\to\ProtoSpace{\sigma},
    \]
    and proto section
    \[
    \PreProtoProtoMin{\sigma}:\ProtoBody{\sigma}\to\PreProtoBody{\sigma}
    \]
    satisfying
    \[
    \PreProtoProj{\sigma}\bigl(\PreProtoBody{\sigma}\bigr)=\ProtoBody{\sigma},
    \qquad
    \PreProtoProj{\sigma}\circ\PreProtoProtoMin{\sigma}
    =
    \mathrm{id}_{\ProtoBody{\sigma}};
    \]
    \item the restriction section
    \[
    \PreProtoSection{\sigma}:\PreProtoSpace{\sigma}\to\SubTarget{\sigma},
    \]
    the induced functional
    \[
    \PreProtoFunctional{\sigma}(q)
    \equiv
    \SubPair{\sigma}\bigl(\PreProtoSection{\sigma}(q),\PreProtoSection{\sigma}(q)\bigr),
    \]
    and the exact shell
    \[
    \PreProtoShellSet{\sigma}
    \equiv
    \bigl(\PreProtoSection{\sigma}\bigr)^{-1}(0)\cap\PreProtoBody{\sigma},
    \]
    with
    \[
    \PreProtoSection{\sigma}\circ\PreProtoProtoMin{\sigma}
    =
    \ProtoSection{\sigma},
    \qquad
    \PreProtoProj{\sigma}\big|_{\PreProtoShellSet{\sigma}}
    :
    \PreProtoShellSet{\sigma}
    \xrightarrow{\cong}
    \ProtoShellSet{\sigma};
    \]
    \item benchmark faithfulness, meaning preservation of the benchmark outputs \eqref{eq:fixedoutputs}, no second operative metric, no enlarged low-energy matter law, and no premature promotion from adoption to completed derivation.
\end{enumerate}
\end{definition}

\begin{definition}[Recorded observer-relevant pre-nucleus exact-shell core]
\label{def:core}
Fix \(\sigma\in\{\Car,\Cur,\Hom\}\). The observer-relevant pre-nucleus exact-shell core \(\PreNucCore\) for sector \(\sigma\) consists of:
\begin{enumerate}
    \item the pre-nucleus state space \(\PreNucSpace{\sigma}\), compact convex body \(\PreNucBody{\sigma}\subset\PreNucSpace{\sigma}\), affine surjection
    \[
    \PreNucProj{\sigma}:\PreNucSpace{\sigma}\to\PreProtoSpace{\sigma},
    \]
    and fiber minimizer
    \[
    \PreNucNucleusMin{\sigma}:\PreProtoBody{\sigma}\to\PreNucBody{\sigma}
    \]
    satisfying
    \[
    \PreNucProj{\sigma}\bigl(\PreNucBody{\sigma}\bigr)=\PreProtoBody{\sigma},
    \qquad
    \PreNucProj{\sigma}\circ\PreNucNucleusMin{\sigma}
    =
    \mathrm{id}_{\PreProtoBody{\sigma}};
    \]
    \item a closed embedded exact pre-nucleus shell
    \[
    \PreNucShellSet{\sigma}\subset\PreNucBody{\sigma}
    \]
    satisfying
    \[
    \PreNucNucleusMin{\sigma}\bigl(\PreProtoShellSet{\sigma}\bigr)
    \subset
    \PreNucShellSet{\sigma},
    \qquad
    \PreNucProj{\sigma}\big|_{\PreNucShellSet{\sigma}}
    :
    \PreNucShellSet{\sigma}
    \xrightarrow{\cong}
    \PreProtoShellSet{\sigma};
    \]
    \item benchmark faithfulness in the sense of Definition \ref{def:nucleus}.
\end{enumerate}
\end{definition}

\begin{definition}[Recorded pre-core exact-shell transport core]
\label{def:transportcore}
Fix \(\sigma\in\{\Car,\Cur,\Hom\}\). The recorded observer-relevant pre-core exact-shell transport core \(\PreCoreTransportCore\) for sector \(\sigma\) consists of:
\begin{enumerate}
    \item a compact transport-core body
    \[
    \PreCoreTrBody{\sigma},
    \]
    an affine projection
    \[
    \PreCoreTrProj{\sigma}:\PreCoreTrBody{\sigma}\to\PreNucBody{\sigma},
    \]
    and a smooth transport-core lift
    \[
    \PreCoreTrLift{\sigma}:\PreNucBody{\sigma}\to\PreCoreTrBody{\sigma}
    \]
    satisfying
    \[
    \PreCoreTrProj{\sigma}\bigl(\PreCoreTrBody{\sigma}\bigr)=\PreNucBody{\sigma},
    \qquad
    \PreCoreTrProj{\sigma}\circ\PreCoreTrLift{\sigma}
    =
    \mathrm{id}_{\PreNucBody{\sigma}};
    \]
    \item a closed embedded exact transport shell
    \[
    \PreCoreTrShell{\sigma}\subset\PreCoreTrBody{\sigma}
    \]
    satisfying
    \[
    \PreCoreTrLift{\sigma}\bigl(\PreNucShellSet{\sigma}\bigr)
    \subset
    \PreCoreTrShell{\sigma},
    \qquad
    \PreCoreTrProj{\sigma}\big|_{\PreCoreTrShell{\sigma}}
    :
    \PreCoreTrShell{\sigma}
    \xrightarrow{\cong}
    \PreNucShellSet{\sigma};
    \]
    \item benchmark faithfulness in the sense of Definition \ref{def:core}.
\end{enumerate}
\end{definition}

\begin{remark}
Definitions \ref{def:fixed}--\ref{def:transportcore} are fixed inputs. The one operative metric, the ordinary matter route, the fixed proto-germ exact-shell target, the recorded current transport nucleus, the recorded pre-nucleus exact-shell core, and the recorded pre-core exact-shell transport core are not reopened here. The present note asks only why the current transport-section core itself should arise from a still more primitive substrate law.
\end{remark}

\section{Exact-Shell-Transport-Section-Generating Substrate Sub-Section Dynamics}

\begin{definition}[Exact-shell-transport-section-generating substrate sub-section dynamics]
\label{def:subsec}
Fix \(\sigma\in\{\Car,\Cur,\Hom\}\). An \emph{exact-shell-transport-section-generating substrate sub-section dynamics package} \(\SubSectionPackage\) for sector \(\sigma\) consists of the following data.
\begin{enumerate}
    \item A finite-dimensional Euclidean sub-section state space \(\SubSectionSpace{\sigma}\), a compact convex sub-section body
    \[
    \SubSectionBody{\sigma}\subset\SubSectionSpace{\sigma}
    \]
    with nonempty interior, and an affine surjection
    \[
    \SubSectionProj{\sigma}:\SubSectionSpace{\sigma}\to\PreCoreTrBody{\sigma}.
    \]
    For each \(w\in\PreCoreTrBody{\sigma}\), write
    \[
    \SubSectionFiber{\sigma}{w}
    \equiv
    \bigl(\SubSectionProj{\sigma}\bigr)^{-1}(w)\cap\SubSectionBody{\sigma}.
    \]
    Each such fiber is required to be nonempty, compact, and convex.
    \item A nonnegative smooth fiber-selection functional
    \[
    \SubSectionFunctional{\sigma}:\SubSectionSpace{\sigma}\to[0,\infty)
    \]
    such that, for every \(w\in\PreCoreTrBody{\sigma}\), the restriction of \(\SubSectionFunctional{\sigma}\) to \(\SubSectionFiber{\sigma}{w}\) is strictly convex and attains a unique minimizer. The resulting minimizer depends smoothly on \(w\).
    \item A closed embedded exact sub-section shell
    \[
    \SubSectionShell{\sigma}\subset\SubSectionBody{\sigma}
    \]
    satisfying
    \[
    \SubSectionShell{\sigma}
    =
    \left\{
    y\in\SubSectionBody{\sigma}:
    \SubSectionProj{\sigma}(y)\in\PreCoreTrShell{\sigma},
    \;
    \SubSectionFunctional{\sigma}(y)
    =
    \min_{z\in\SubSectionFiber{\sigma}{\SubSectionProj{\sigma}(y)}}
    \SubSectionFunctional{\sigma}(z)
    \right\},
    \]
    and
    \[
    \SubSectionProj{\sigma}\big|_{\SubSectionShell{\sigma}}
    :
    \SubSectionShell{\sigma}
    \xrightarrow{\cong}
    \PreCoreTrShell{\sigma}
    \]
    is a diffeomorphism.
    \item Benchmark faithfulness, meaning preservation of the benchmark outputs \eqref{eq:fixedoutputs}, no second operative metric, no enlarged low-energy matter law, and no premature promotion from adoption to completed derivation.
\end{enumerate}
\end{definition}

\begin{remark}
Definition \ref{def:subsec} is intentionally written over the already fixed pre-core exact-shell transport body and shell. It does not restore the older image-core or full sub-pre-core presentation as the live burden. Its positive role is narrower: the current transport section is generated as the unique fiberwise minimizer of a deeper substrate selector law.
\end{remark}

\section{Induced Observer-Relevant Sub-Pre-Core Exact-Shell Transport-Section Core}

\begin{definition}[Observer-relevant sub-pre-core exact-shell transport-section core]
\label{def:sectioncore}
Fix \(\sigma\in\{\Car,\Cur,\Hom\}\) and a benchmark-faithful sub-section dynamics package in the sense of Definition \ref{def:subsec}. The \emph{observer-relevant sub-pre-core exact-shell transport-section core} \(\SubPreSectionCore\) for sector \(\sigma\) consists of:
\begin{enumerate}
    \item the canonical transport section
    \[
    \SubPreSection{\sigma}(w)
    \equiv
    \operatorname*{arg\,min}_{y\in\SubSectionFiber{\sigma}{w}}
    \SubSectionFunctional{\sigma}(y),
    \qquad
    w\in\PreCoreTrBody{\sigma};
    \]
    \item benchmark faithfulness in the sense of Definition \ref{def:subsec}.
\end{enumerate}
The shell restriction
\[
\SubPreSection{\sigma}\big|_{\PreCoreTrShell{\sigma}}
:
\PreCoreTrShell{\sigma}\to\SubSectionShell{\sigma}
\]
is understood as derived data rather than as separately primitive core data.
\end{definition}

\begin{remark}
Definition \ref{def:sectioncore} is exactly the current live benchmark-facing datum, now written as a canonical consequence of the deeper sub-section dynamics rather than as a remaining primitive stopping point.
\end{remark}

\section{Canonical Recovery of the Current Section Core}

\begin{proposition}[The canonical transport section is well defined]
\label{prop:section}
For every benchmark-faithful sub-section dynamics package, the map \(\SubPreSection{\sigma}\) of Definition \ref{def:sectioncore} is a smooth section of \(\SubSectionProj{\sigma}\), so that
\[
\SubSectionProj{\sigma}\circ\SubPreSection{\sigma}
=
\mathrm{id}_{\PreCoreTrBody{\sigma}}.
\]
\end{proposition}

\begin{proof}
Definition \ref{def:subsec}(1) makes each fiber \(\SubSectionFiber{\sigma}{w}\) nonempty, compact, and convex. Definition \ref{def:subsec}(2) states that \(\SubSectionFunctional{\sigma}\) has a unique minimizer on each such fiber and that the minimizer depends smoothly on \(w\). Definition \ref{def:sectioncore}(1) therefore gives a well-defined smooth map
\[
\SubPreSection{\sigma}:\PreCoreTrBody{\sigma}\to\SubSectionBody{\sigma}.
\]
Because every minimizer lies in the fiber over the point at which it is selected, we have
\[
\SubSectionProj{\sigma}\bigl(\SubPreSection{\sigma}(w)\bigr)=w
\]
for every \(w\in\PreCoreTrBody{\sigma}\), which is the displayed identity.
\end{proof}

\begin{proposition}[The exact shell is the shell restriction of the canonical section]
\label{prop:shell}
For every benchmark-faithful sub-section dynamics package,
\[
\SubSectionShell{\sigma}
=
\SubPreSection{\sigma}\bigl(\PreCoreTrShell{\sigma}\bigr),
\]
and the shell restriction
\[
\SubPreSection{\sigma}\big|_{\PreCoreTrShell{\sigma}}
:
\PreCoreTrShell{\sigma}\xrightarrow{\cong}\SubSectionShell{\sigma}
\]
is a diffeomorphism with inverse
\[
\SubSectionProj{\sigma}\big|_{\SubSectionShell{\sigma}}
:
\SubSectionShell{\sigma}\xrightarrow{\cong}\PreCoreTrShell{\sigma}.
\]
\end{proposition}

\begin{proof}
Let \(w\in\PreCoreTrShell{\sigma}\). By Definition \ref{def:sectioncore}(1), \(\SubPreSection{\sigma}(w)\) is the unique minimizer of \(\SubSectionFunctional{\sigma}\) on the fiber \(\SubSectionFiber{\sigma}{w}\). Since \(w\in\PreCoreTrShell{\sigma}\), Definition \ref{def:subsec}(3) implies
\[
\SubPreSection{\sigma}(w)\in\SubSectionShell{\sigma}.
\]
Hence
\[
\SubPreSection{\sigma}\bigl(\PreCoreTrShell{\sigma}\bigr)\subset\SubSectionShell{\sigma}.
\]

Conversely, if \(y\in\SubSectionShell{\sigma}\), then \(w\equiv\SubSectionProj{\sigma}(y)\) lies in \(\PreCoreTrShell{\sigma}\), and Definition \ref{def:subsec}(3) says that \(y\) is a minimizer of \(\SubSectionFunctional{\sigma}\) on \(\SubSectionFiber{\sigma}{w}\). By uniqueness of that minimizer, \(y=\SubPreSection{\sigma}(w)\). This proves the displayed equality.

The map \(\SubSectionProj{\sigma}\big|_{\SubSectionShell{\sigma}}\) is a diffeomorphism by Definition \ref{def:subsec}(3). The displayed equality shows that its inverse is exactly \(\SubPreSection{\sigma}\big|_{\PreCoreTrShell{\sigma}}\).
\end{proof}

\begin{proposition}[The induced transport section is canonical]
\label{prop:canonical}
Fix a benchmark-faithful sub-section dynamics package. If
\[
s:\PreCoreTrBody{\sigma}\to\SubSectionBody{\sigma}
\]
is any smooth section satisfying
\[
\SubSectionProj{\sigma}\circ s
=
\mathrm{id}_{\PreCoreTrBody{\sigma}}
\]
and
\[
\SubSectionFunctional{\sigma}\bigl(s(w)\bigr)
=
\min_{y\in\SubSectionFiber{\sigma}{w}}
\SubSectionFunctional{\sigma}(y)
\qquad
\text{for every }
w\in\PreCoreTrBody{\sigma},
\]
then \(s=\SubPreSection{\sigma}\). In particular, the induced observer-relevant section core carries no residual choice once the deeper sub-section dynamics are fixed.
\end{proposition}

\begin{proof}
For each \(w\in\PreCoreTrBody{\sigma}\), the point \(s(w)\) belongs to \(\SubSectionFiber{\sigma}{w}\) and minimizes \(\SubSectionFunctional{\sigma}\) there. By Definition \ref{def:subsec}(2), that minimizer is unique. Therefore
\[
s(w)=\SubPreSection{\sigma}(w)
\]
for every \(w\), hence \(s=\SubPreSection{\sigma}\).
\end{proof}

\begin{proposition}[Downstream exact-shell transport factors through the induced section core]
\label{prop:downstream}
For every benchmark-faithful sub-section dynamics package, all downstream shell transport to the recorded pre-nucleus exact-shell core, the current transport nucleus, and the fixed proto-germ exact-shell target already factors through the induced section core \(\SubPreSectionCore\). More precisely,
\[
\PreCoreTrProj{\sigma}\circ
\SubSectionProj{\sigma}\big|_{\SubSectionShell{\sigma}}
:
\SubSectionShell{\sigma}\xrightarrow{\cong}\PreNucShellSet{\sigma},
\]
\[
\PreNucProj{\sigma}\circ\PreCoreTrProj{\sigma}\circ
\SubSectionProj{\sigma}\big|_{\SubSectionShell{\sigma}}
:
\SubSectionShell{\sigma}\xrightarrow{\cong}\PreProtoShellSet{\sigma},
\]
and
\[
\PreProtoProj{\sigma}\circ\PreNucProj{\sigma}\circ\PreCoreTrProj{\sigma}\circ
\SubSectionProj{\sigma}\big|_{\SubSectionShell{\sigma}}
:
\SubSectionShell{\sigma}\xrightarrow{\cong}\ProtoShellSet{\sigma}
\]
are diffeomorphisms.
\end{proposition}

\begin{proof}
By Proposition \ref{prop:shell},
\[
\SubSectionProj{\sigma}\big|_{\SubSectionShell{\sigma}}
:
\SubSectionShell{\sigma}\xrightarrow{\cong}\PreCoreTrShell{\sigma}
\]
is a diffeomorphism. Definition \ref{def:transportcore}(2) gives a diffeomorphism
\[
\PreCoreTrProj{\sigma}\big|_{\PreCoreTrShell{\sigma}}
:
\PreCoreTrShell{\sigma}\xrightarrow{\cong}\PreNucShellSet{\sigma}.
\]
Their composition is the first displayed diffeomorphism.

Definition \ref{def:core}(2) gives
\[
\PreNucProj{\sigma}\big|_{\PreNucShellSet{\sigma}}
:
\PreNucShellSet{\sigma}\xrightarrow{\cong}\PreProtoShellSet{\sigma},
\]
so composing with the first diffeomorphism gives the second displayed map. Definition \ref{def:nucleus}(2) then gives
\[
\PreProtoProj{\sigma}\big|_{\PreProtoShellSet{\sigma}}
:
\PreProtoShellSet{\sigma}\xrightarrow{\cong}\ProtoShellSet{\sigma},
\]
and composing once more yields the third displayed diffeomorphism. Thus nothing benchmark-facing downstream of the current transport section requires more than the induced section core together with the already fixed transport-core data.
\end{proof}

\begin{proposition}[The benchmark package remains fixed]
\label{prop:benchmark}
Every benchmark-faithful exact-shell-transport-section-generating substrate sub-section dynamics package preserves the same benchmark one-metric package \eqref{eq:fixedoutputs}. In particular, the operative metric remains exactly \(\CompletedMetric\), the ordinary matter route is unchanged, there is no second operative metric, and the low-energy matter law is not enlarged.
\end{proposition}

\begin{proof}
This is part of benchmark faithfulness in Definition \ref{def:subsec}(4). The present note only derives the current section core from deeper substrate selector data. It does not alter the benchmark exact-closure package.
\end{proof}

\section{Main Theorem}

\begin{theorem}[Observer-relevant sub-pre-core exact-shell transport-section core from exact-shell-transport-section-generating substrate sub-section dynamics]
\label{thm:main}
Fix the primitive continuation data of Definition \ref{def:fixed}, the recorded proto-germ exact-shell target of Definition \ref{def:proto}, the recorded current transport nucleus of Definition \ref{def:nucleus}, the recorded observer-relevant pre-nucleus exact-shell core of Definition \ref{def:core}, the recorded observer-relevant pre-core exact-shell transport core of Definition \ref{def:transportcore}, and a benchmark-faithful exact-shell-transport-section-generating substrate sub-section dynamics package in the sense of Definition \ref{def:subsec}. Then, for each sector \(\sigma\in\{\Car,\Cur,\Hom\}\), the current observer-relevant sub-pre-core exact-shell transport-section core is no longer an unexplained primitive stopping point on the active line: it is recovered canonically from the deeper sub-section dynamics. More precisely:
\begin{enumerate}
    \item the deeper sub-section dynamics generate the observer-relevant sub-pre-core exact-shell transport-section core \(\SubPreSectionCore\) by Proposition \ref{prop:section};
    \item the exact section shell is exactly the shell restriction of that canonical transport section by Proposition \ref{prop:shell};
    \item the induced transport section is canonical and carries no residual minimizing-choice freedom once the deeper package is fixed by Proposition \ref{prop:canonical};
    \item all downstream shell transport to the recorded pre-nucleus exact-shell core, the current transport nucleus, and the fixed proto-germ exact-shell target factors through the induced section core by Proposition \ref{prop:downstream};
    \item the benchmark boundary remains fixed by Proposition \ref{prop:benchmark};
    \item consequently the remaining honest burden below the previous section-core frontier is no longer the observer-relevant sub-pre-core exact-shell transport-section core \(\SubPreSectionCore\) as such. What remains is to derive or justify the deeper exact-shell-transport-section-generating substrate sub-section dynamics \(\SubSectionPackage\) from still more primitive substrate structure, or else prove a still sharper theorem-level collapse strictly inside that deeper package.
\end{enumerate}
\end{theorem}

\begin{proof}
Item (1) is Proposition \ref{prop:section}. Item (2) is Proposition \ref{prop:shell}. Item (3) is Proposition \ref{prop:canonical}. Item (4) is Proposition \ref{prop:downstream}. Item (5) is Proposition \ref{prop:benchmark}. Item (6) is the routing consequence of the previous items together with the active safeguard against counting another same-output deeper-origin relay as new front-line progress when the live benchmark-relevant datum is unchanged.
\end{proof}

\begin{corollary}[What must be done next]
\label{cor:next}
After Theorem \ref{thm:main}, the next honest benchmark-facing manuscript must either:
\begin{enumerate}
    \item derive or justify the exact-shell-transport-section-generating substrate sub-section dynamics \(\SubSectionPackage\) from still more primitive substrate structure in a way that changes benchmark-facing status;
    \item identify a genuinely smaller theorem-level observer-relevant datum strictly inside \(\SubSectionPackage\) and prove a sharper collapse to it;
    \item do either of the previous two while preserving the same benchmark one-metric package and the same claim boundary.
\end{enumerate}
Another same-output deeper-origin relay that merely preserves the generated section core, another re-expansion back to the larger transport-image or full sub-pre-core presentations as if they were still independently live, another reopening of the already-screened larger pre-core packages as if they were again the active burden, or any move that introduces a second operative metric, enlarges the ordinary matter law, or uses stronger promotion language does not satisfy the present burden. If a quantum continuation is pursued, it matters benchmark-facingly only insofar as it derives this same deeper sub-section dynamics, collapses them further, or closes a benchmark derivation gate in a genuinely new way.
\end{corollary}

\begin{proof}
Theorem \ref{thm:main} shows that the current section core is no longer the deepest unexplained datum on the active line. Therefore a subsequent manuscript must improve on the deeper sub-section dynamics rather than merely restate the already generated section core.
\end{proof}

\section{Conclusion}

The positive result of this note is that the current observer-relevant sub-pre-core exact-shell transport-section core is no longer primitive on the active derivational line. Once one works with an exact-shell-transport-section-generating substrate sub-section dynamics package over the already fixed pre-core transport body, the section is forced as the unique fiberwise minimizer of a deeper substrate selection law. The exact shell of that package is exactly the shell restriction of the same section, and all downstream shell transport to the recorded pre-nucleus exact-shell core, the current transport nucleus, and the fixed proto-germ exact-shell target already factors through that induced datum.

This preserves the same benchmark one-metric package and the same claim boundary. The result does not yet provide a completed first-principles derivation of GR from substrate microphysics, and it does not license stronger promotion language. Its narrower benchmark-facing gain is that the remaining honest burden now sits one level lower again: not on the current section core itself, but on the deeper sub-section dynamics that canonically generate it, or on a still smaller collapse strictly inside that deeper package.

\end{document}
